\chapter{Analisi dei requisiti}

La struttura di questo capitolo si ispira allo standard ISO/IEC/IEEE 29148:2018 - Systems and Software Engineering - Life cycle processes- Requirements engineering
\textbf{[CITAZIONE: ISO/IEC/IEEE 29148:2018]}
DA CAPIRE CITAZIONE.

% -----------------------------------

\section{Requisiti funzionali}

\subsection{Acquisizione dell'input}
TODO: testo da scrivere.

\subsection{Estrazione del testo}
TODO: testo da scrivere.

\subsection{Rilevamento delle PII}
TODO: testo da scrivere.

\subsection{Aggregazione delle PII frammentate | DA VALUTARE}
TODO: testo da scrivere.

\subsection{Costruzione di una knowledge base delle entità | DA VALUTARE}
TODO: testo da scrivere.

\subsection{Ingestionde del/i ROPA}
TODO: testo da scrivere.

\subsection{Confronto knowledge (o risultati di analisi) con ROPA}
TODO: testo da scrivere.

\subsection{Rilevamento delle anomalie di conformità}
TODO: testo da scrivere.

\subsection{Generazione report per DPO}
TODO: testo da scrivere.

% -----------------------------------

\section{Requisiti normativi}

\subsection{Conformità ai principi di limitazione delle finalità e della conservazione}
TODO: testo da scrivere.

\subsection{Riferimento al registro delle attività di trattamento come fonte normativa}
TODO: testo da scrivere.

\subsection{Copertura delle categorie di dati degni di particolare attenzione}
TODO: testo da scrivere.

\subsection{Posizionamento rispetto all'art.32(1)(d) GDPR}
TODO: testo da scrivere.

\subsection{Esclusione delle basi giuridiche dal perimetro di analisi automatica}
TODO: testo da scrivere.

% -----------------------------------

\section{Requisiti operativi}

\subsection{Elaborazione locale dei dati}
TODO: testo da scrivere.

\subsection{Vincolo sui modelli AI impiegabili}
TODO: testo da scrivere.

\subsection{Facilità installazione e aggiornamento}
TODO: testo da scrivere.

\subsection{Estensibilità delle categorie e detectors}
TODO: testo da scrivere.

\subsection{Minimizzazione dei dati trattati dal sistema stesso}
TODO: testo da scrivere.

% -----------------------------------

\section{Requisiti di qualità}
TODO: testo da scrivere.

\subsection{Prestazioni su larga scala}
TODO: testo da scrivere.

\subsection{Precisione del rilevamento}
TODO: testo da scrivere.

\subsection{Tracciabilità delle anomalie}
TODO: testo da scrivere.

% -----------------------------------

\section{Requisiti di interfaccia}

\subsection{\textbf{da definire}}
TODO: testo da scrivere.
